\documentclass[11pt,a4paper]{article}
\usepackage[margin=1in]{geometry}
\usepackage{amsmath,amssymb,amsthm}
\usepackage{algorithm}
\usepackage{algpseudocode}
\usepackage{booktabs}
\usepackage{hyperref}

\theoremstyle{definition}
\newtheorem{definition}{Definition}[section]
\newtheorem{theorem}{Theorem}[section]
\newtheorem{lemma}[theorem]{Lemma}
\newtheorem{remark}{Remark}[section]

\title{\textbf{Rigorous Mathematical Specification, Period Dynamics, and Convergence of Cosine Annealing and SGDR Warm Restarts}}
\author{Team Scaibu \\ \small Email: \texttt{jaydeep.9664920749@gmail.com} \\ \small Architecture and Core Engine Team}
\date{October 2026}

\begin{document}
\maketitle

\begin{abstract}
This document presents the complete theoretical formulation and algorithmic specification of Cosine Annealing and Stochastic Gradient Descent with Warm Restarts (SGDR). While step-decay schedules require arbitrary manual epoch milestones, cosine annealing decreases learning rates along a continuous half-cosine trajectory. SGDR introduces periodic warm restarts where the learning rate resets to its peak, encouraging the optimizer to escape suboptimal local basins and traverse multi-modal loss landscapes. We formalize geometric period scaling ($T_{\text{mult}}$), present the exact algorithmic pseudo-code matching the reference implementation, derive convergence bounds, compute a worked numerical example, and discuss numerical implementation considerations.
\end{abstract}

\section{Notation and Mathematical Preliminaries}

Let $t \in \mathbb{N}_{\ge 0}$ denote the discrete global training step.

\begin{itemize}
    \item $t$: Current training iteration ($t \ge 0$).
    \item $T_{\text{max}} \equiv T_0 \in \mathbb{N}_{\ge 1}$: Base period duration (steps in initial cycle).
    \item $\eta_{\max} > 0$: Maximum initial learning rate.
    \item $\eta_{\min} \ge 0$: Lower bound learning rate floor (default $0.0$).
    \item $T_{\text{mult}} \ge 1$: Multiplicative factor expanding cycle period for subsequent restart epochs ($T_i = T_0 \cdot T_{\text{mult}}^i$).
    \item $i \in \mathbb{N}_{\ge 0}$: Current restart cycle index.
    \item $T_{\text{cur}}$: Relative step index elapsed within the current cycle $i$ ($0 \le T_{\text{cur}} < T_i$).
\end{itemize}

\begin{table}[h]
\centering
\caption{Summary of Mathematical Symbols for Cosine Decay and SGDR}
\begin{tabular}{lll}
\toprule
\textbf{Symbol} & \textbf{Type / Domain} & \textbf{Description} \\
\midrule
$t$ & $\mathbb{N}_{\ge 0}$ & Global step counter \\
$T_0$ & $\mathbb{N}_{\ge 1}$ & Base period / total steps \\
$\eta_{\max}$ & $\mathbb{R}_{> 0}$ & Peak learning rate \\
$\eta_{\min}$ & $\mathbb{R}_{\ge 0}$ & Minimum learning rate \\
$T_{\text{mult}}$ & $\mathbb{N}_{\ge 1}$ & Period expansion factor \\
$i$ & $\mathbb{N}_{\ge 0}$ & Cycle index \\
$T_{\text{cur}}$ & $[0, T_i)$ & Steps elapsed in cycle $i$ \\
$\eta_t$ & $[\eta_{\min}, \eta_{\max}]$ & Effective scheduled learning rate \\
\bottomrule
\end{tabular}
\end{table}

\section{Problem Definition and Core Concepts}

\subsection{Intuition and Motivation}
In non-convex neural optimization, static and monotonically decaying learning rates risk settling in sharp local minima or flat saddle regions early in training. Cosine annealing smoothly decays the learning rate according to a trigonometric curve:
\begin{equation}
\eta(t) = \eta_{\min} + \frac{1}{2}(\eta_{\max} - \eta_{\min}) \left(1 + \cos\left(\frac{\pi t}{T}\right)\right)
\end{equation}
When SGDR warm restarts are enabled, every time $T_{\text{cur}}$ completes a cycle $T_i$, $\eta$ instantaneously resets to $\eta_{\max}$. This sudden surge in kinetic energy enables the optimization trajectory to break free from shallow basins and discover wider, better-generalizing minima across multiple cycles (often ensembled via Snapshot Ensembling).

\subsection{Formal Mathematical Definition}
\begin{definition}[SGDR Periodic Formulation]
Let $T_i = T_0 \cdot T_{\text{mult}}^i$ denote the duration of the $i$-th cycle. The cumulative steps before cycle $i$ is:
\begin{equation}
S_i = \sum_{k=0}^{i-1} T_k = T_0 \sum_{k=0}^{i-1} T_{\text{mult}}^k = T_0 \frac{T_{\text{mult}}^i - 1}{T_{\text{mult}} - 1} \quad (\text{for } T_{\text{mult}} > 1)
\end{equation}
Given global step $t$, the current cycle $i$ is the unique integer satisfying $S_i \le t < S_{i+1}$, and the relative progress within cycle $i$ is $T_{\text{cur}} = t - S_i$. The learning rate is:
\begin{equation}
\eta_t = \eta_{\min} + \frac{1}{2} (\eta_{\max} - \eta_{\min}) \left( 1 + \cos\left( \frac{\pi T_{\text{cur}}}{T_i} \right) \right) \label{eq:sgdr_formula}
\end{equation}
\end{definition}

\section{Algorithmic Specification}

The computational pipeline implemented in \texttt{NnAlgoCosineDecayRestarts} is shown below.

\begin{algorithm}[H]
\caption{Cosine Annealing and SGDR Warm Restarts Schedule}
\begin{algorithmic}[1]
\Require Current step $t \ge 0$, base steps $T_0 \ge 1$, peak rate $\eta_{\max} > 0$
\Require Minimum rate $\eta_{\min} \ge 0$, flag $\text{use\_restarts} \in \{\textbf{true}, \textbf{false}\}$, multiplier $T_{\text{mult}} \ge 1$
\Ensure Effective learning rate $\eta_t$, cycle index $\text{current\_cycle}$
\If{$t < 0$ \textbf{or} $T_0 < 1$ \textbf{or} $\eta_{\max} \le 0$ \textbf{or} $\eta_{\min} < 0$ \textbf{or} $T_{\text{mult}} < 1$}
    \State \textbf{throw} PreconditionFailureException(``Invalid schedule arguments'')
\EndIf
\If{\textbf{not} $\text{use\_restarts}$}
    \State $t_{\text{clamped}} \gets \min(t, T_0)$
    \State $p \gets \text{float}(t_{\text{clamped}}) / \text{float}(T_0)$
    \State $\eta_t \gets \eta_{\min} + 0.5 \cdot (\eta_{\max} - \eta_{\min}) \cdot (1.0 + \cos(\pi \cdot p))$
    \State \Return $\{\text{learning\_rate}: \eta_t, \text{current\_cycle}: 0\}$
\EndIf
\State $T_{\text{cur}} \gets t, \quad T_i \gets T_0, \quad \text{cycle} \gets 0$
\While{$T_{\text{cur}} \ge T_i$}
    \State $T_{\text{cur}} \gets T_{\text{cur}} - T_i$
    \State $T_i \gets T_i \cdot T_{\text{mult}}$
    \State $\text{cycle} \gets \text{cycle} + 1$
\EndWhile
\State $p \gets \text{float}(T_{\text{cur}}) / \text{float}(T_i)$
\State $\eta_t \gets \eta_{\min} + 0.5 \cdot (\eta_{\max} - \eta_{\min}) \cdot (1.0 + \cos(\pi \cdot p))$
\State \Return $\{\text{learning\_rate}: \eta_t, \text{current\_cycle}: \text{cycle}\}$
\end{algorithmic}
\end{algorithm}

\section{Theoretical Guarantees and Annealing Rates}

\begin{theorem}[Continuity and Vanishing Derivative at Boundary]
The cosine schedule $\eta(p) = \eta_{\min} + \frac{1}{2}(\eta_{\max} - \eta_{\min})(1 + \cos(\pi p))$ is $\mathcal{C}^\infty$ smooth on $(0, 1)$ with zero first derivative at both boundary extremes:
\begin{equation}
\left.\frac{d\eta}{dp}\right|_{p=0} = 0 \quad \text{and} \quad \left.\frac{d\eta}{dp}\right|_{p=1} = 0
\end{equation}
\end{theorem}

\begin{proof}
Differentiating with respect to progress $p$:
\begin{equation}
\frac{d\eta}{dp} = -\frac{\pi}{2} (\eta_{\max} - \eta_{\min}) \sin(\pi p)
\end{equation}
Evaluating at $p=0$: $\sin(0) = 0 \implies \frac{d\eta}{dp} = 0$. Evaluating at $p=1$: $\sin(\pi) = 0 \implies \frac{d\eta}{dp} = 0$.
The zero derivative at $p=0$ ensures smooth transition from warmup/peak rate without abrupt shock, and zero derivative at $p=1$ stabilizes fine convergence at the bottom of the basin.
\end{proof}

\subsection{Limitations}
\begin{itemize}
    \item \textbf{Total Horizon Sensitivity}: Without restarts, if training is extended past $T_{\text{max}}$, $\eta_t$ stays frozen at $\eta_{\min}$, requiring manual schedule re-anchoring.
\end{itemize}

\section{Complexity Analysis}

\subsection{Time and Space Complexity}
\begin{itemize}
    \item \textbf{Time}: While loop executes at most $\mathcal{O}(\log_{T_{\text{mult}}} (t / T_0))$ iterations (typically $\le 5$ cycles total), followed by 1 cosine evaluation ($\mathcal{O}(1)$ time).
    \item \textbf{Space}: $\mathcal{O}(1)$ memory.
\end{itemize}

\section{Worked Numerical Example}

Consider $T_0 = 100$, $T_{\text{mult}} = 2$, $\eta_{\max} = 0.1$, $\eta_{\min} = 0.0$, $\text{use\_restarts} = \textbf{true}$, at step $t = 150$.

\begin{enumerate}
    \item \textbf{Cycle 0}: $T_0 = 100$. Since $t=150 \ge 100$:
    \begin{equation}
    T_{\text{cur}} = 150 - 100 = 50, \quad T_1 = 100 \times 2 = 200, \quad \text{cycle} = 1
    \end{equation}
    \item Since $T_{\text{cur}} = 50 < T_1 = 200$, the while loop terminates with $\text{current\_cycle} = 1$.
    \item \textbf{Progress in Cycle 1}:
    \begin{equation}
    p = \frac{50}{200} = 0.25
    \end{equation}
    \item \textbf{Learning Rate}:
    \begin{align}
    \eta_{150} &= 0.0 + 0.5(0.1 - 0.0)(1.0 + \cos(0.25 \pi)) \\
    &= 0.05 \times (1.0 + 0.70710678) = 0.05 \times 1.70710678 \approx 0.0853553
    \end{align}
\end{enumerate}

\section{Numerical Considerations and Edge Cases}

\begin{itemize}
    \item \textbf{Step Clamp Without Restarts}: When $\text{use\_restarts} = \textbf{false}$ and $t > T_0$, clamping $t_{\text{clamped}} = T_0$ ensures $p=1.0$ and $\cos(\pi) = -1.0$, stably holding $\eta_t = \eta_{\min}$.
    \item \textbf{Floating-Point Trigonometric Accuracy}: Using high-precision $\pi \approx 3.141592653589793$ avoids phase drift in cosine calculations over millions of steps.
\end{itemize}

\begin{thebibliography}{9}
\bibitem{loshchilov2017sgdr} I.~Loshchilov and F.~Hutter, ``SGDR: Stochastic gradient descent with warm restarts,'' in \emph{International Conference on Learning Representations (ICLR)}, 2017.
\bibitem{huang2017} G.~Huang et al., ``Snapshot ensembles: Train 1, get m for free,'' in \emph{International Conference on Learning Representations (ICLR)}, 2017.
\bibitem{gotmare2018} A.~Gotmare et al., ``A closer look at deep learning heuristics: Learning rate restarts, warmup and distill,'' in \emph{International Conference on Learning Representations (ICLR)}, 2019.
\end{thebibliography}

\end{document}
